\documentclass[11pt,a4paper,spanish]{article}

\usepackage[utf-8]{inputenc}
\usepackage[spanish]{babel}
\usepackage{xcolor}
\usepackage{amssymb}
\usepackage{listings}
\usepackage{fancyhdr}
\usepackage{geometry}
\usepackage{hyperref}
\usepackage{array}
\usepackage{tabulary}
\usepackage{booktabs}
\usepackage{tikz}
\usetikzlibrary{babel,arrows.meta,positioning,calc,fit,backgrounds,shapes.misc}

\geometry{margin=0.75in}
\pagestyle{fancy}
\fancyhf{}
\rhead{\small Hermes Stack}
\lhead{\small Manual de Uso}
\cfoot{\thepage}

\lstset{
  basicstyle=\ttfamily\footnotesize,
  breaklines=true,
  columns=fullflexible,
  language=bash,
  keywordstyle=\color{blue}\bfseries,
  commentstyle=\color{gray},
  stringstyle=\color{red},
  showstringspaces=false,
  frame=single,
  rulecolor=\color{gray!30},
  aboveskip=6pt,
  belowskip=6pt
}

\definecolor{hermes-dark}{RGB}{45, 45, 45}
\definecolor{hermes-blue}{RGB}{68, 134, 232}
\definecolor{hermes-green}{RGB}{52, 214, 146}

\title{\huge\textbf{Manual de Uso}\\[\smallskip]\Large Claude Code \textbullet{} Hermes Stack\\[\smallskip]\normalsize Tu referencia rápida. Diseñado para escanear, no para leer.}
\author{}
\date{\small Actualizado: 25 de mayo de 2026}

\begin{document}

\maketitle

\section*{⚡ Panel de Control — Respuesta en < 5 segundos}

Cuando necesites algo, aquí está la acción correcta:

\vspace{0.3cm}

\noindent\begin{tabular}{|p{2.5cm}|p{3.5cm}|p{3.5cm}|}
\hline
\textbf{Situación} & \textbf{Comando / Acción} & \textbf{Por qué} \\
\hline
Inicio de sesión & \texttt{make doctor} & Estado completo del stack en 1 comando \\
¿Qué tengo pendiente? & \texttt{cat /root/PENDIENTES.md | head -30} & Tu backlog ordenado por prioridad \\
El stack no responde & \texttt{/pipeline} & Diagnóstico en lenguaje simple + guía \\
Quiero cambiar código & Pídelo directamente en el chat & Claude lo hace — no necesitas comandos \\
¿Funcionó el cambio? & \texttt{make check} & Build + health + lint en 1 comando \\
Quiero hacer un commit & Pídele a Claude que lo haga & Usa el patrón HTTPS correcto automáticamente \\
Sesión larga, quiero parar & Pide ``haz clock-out'' & Claude actualiza PENDIENTES, commit, handoff \\
Nueva sesión, ¿dónde estaba? & \texttt{cat /root/SESSION\_HANDOFF.md} & El paso siguiente está escrito literalmente \\
Quiero que Claude recuerde & ``recuerda que...'' en el chat & Lo guarda permanentemente en memoria \\
El harness está desactualizado & \texttt{/evolve} & Audita y actualiza CLAUDE.md automáticamente \\
\hline
\end{tabular}

\section*{🔄 Flujo de Sesión}

Una sesión tiene 5 fases. Siempre en este orden:

\vspace{0.4cm}

\subsection*{Fase 1: Clock-In}
Verificar estado real del stack y qué hay pendiente:
\begin{lstlisting}
make doctor                    # diagnóstico completo
cat /root/PENDIENTES.md        # ver tareas activas
cat /root/SESSION_HANDOFF.md   # retomar desde donde terminó
\end{lstlisting}

\subsection*{Fase 2: Elegir Tarea (UNA sola)}
\begin{itemize}
  \item 🔴 Prioridad alta primero (items hs-001, hs-002)
  \item Si surge idea secundaria → registrar en PENDIENTES.md pero \textbf{no empezarla}
  \item Una tarea a la vez previene el ``70\% en 3 tareas''
\end{itemize}

\subsection*{Fase 3: Ejecutar (en el chat)}
\begin{itemize}
  \item Describir qué quieres → Claude investiga + implementa
  \item Si es algo riesgoso → Claude te pide confirmación
  \item Si usas voz → dicta directo, Claude normaliza errores tipográficos
\end{itemize}

\subsection*{Fase 4: Verificar (Definition of Done — 4 criterios)}
Una tarea está done cuando se cumplen TODOS estos:
\begin{enumerate}
  \item \checkmark Código implementado
  \item \checkmark \texttt{make check} → sin errores
  \item \checkmark Evidencia en PENDIENTES.md (commit hash o output real)
  \item \checkmark \texttt{docker compose ps} → todos los servicios \texttt{Up}
\end{enumerate}

\subsection*{Fase 5: Clock-Out}
Al terminar sesión:
\begin{lstlisting}
# Pedir: "haz clock-out"
# Claude hará:
#  1. Actualizar PENDIENTES.md con evidencia
#  2. Ejecutar make check para verificar stack sano
#  3. Hacer commit de todo el trabajo
#  4. Crear SESSION_HANDOFF.md
\end{lstlisting}

\vspace{0.5cm}

\section{🎯 Slash Commands}

Comandos especializados que activan comportamientos automáticos:

\vspace{0.3cm}

\begin{tabular}{|p{2cm}|p{3cm}|p{4cm}|}
\hline
\textbf{Comando} & \textbf{Cuándo usarlo} & \textbf{Qué hace} \\
\hline
\texttt{/evolve} & Al final de sesiones largas o si algo no está documentado & Audita los 5 subsistemas del harness, extrae aprendizajes, actualiza CLAUDE.md \\
\texttt{/pipeline} & El stack falla y no sabes por dónde empezar & Ejecuta \texttt{make doctor} y explica en lenguaje simple qué está mal \\
\hline
\end{tabular}

\vspace{0.3cm}

\noindent\textbf{Tip para ADHD:} Estos comandos son de mantenimiento. El chat normal cubre el 95\% de lo que necesitas.

\section{🔧 Skills Ruflo — Cuándo Activarlos}

Escribe la frase clave en el chat y Claude usa el skill automáticamente:

\subsection*{Diagnóstico y Estado}
\begin{itemize}
  \item ``estado del stack'' / ``qué está corriendo'' → \texttt{ruflo-status}
  \item ``diagnóstico completo'' → \texttt{ruflo-doctor}
  \item ``audita la seguridad'' → \texttt{security-audit}
\end{itemize}

\subsection*{Memoria y Contexto}
\begin{itemize}
  \item ``busca en tu memoria...'' / ``recuerdas cuando...'' → \texttt{recall}
  \item ``guarda esto en memoria'' / ``recuerda que...'' → \texttt{ruflo-memory}
\end{itemize}

\subsection*{Código y Documentación}
\begin{itemize}
  \item ``genera tests para [archivo]'' → \texttt{testgen}
  \item ``documenta [módulo]'' → \texttt{ruflo-docs}
  \item ``revisa este PR'' / ``revisa el diff'' → \texttt{code-review}
\end{itemize}

\subsection*{Automatización y Background}
\begin{itemize}
  \item ``programa X cada [intervalo]'' → \texttt{ruflo-loop}
  \item ``activa autopilot para...'' → \texttt{autopilot}
\end{itemize}

\subsection*{Rastreo de Costos}
\begin{itemize}
  \item ``cuánto costó esta sesión'' → \texttt{ruflo-cost}
  \item ``optimiza mi uso de tokens'' → \texttt{cost-optimize}
\end{itemize}

\section{🌳 Árbol de Decisión}

\vspace{0.2cm}

\noindent¿Qué necesitas ahora?

\begin{itemize}
  \item \textbf{Stack no responde / algo está caído}
    \begin{itemize}
      \item \texttt{/pipeline} → si no basta: \texttt{make logs} → \texttt{docker compose ps}
    \end{itemize}
  \item \textbf{Quiero cambiar código / funcionalidad}
    \begin{itemize}
      \item Describir en el chat (Claude implementa) → \texttt{make check} → pedir commit
    \end{itemize}
  \item \textbf{¿Qué hay pendiente?}
    \begin{itemize}
      \item \texttt{cat /root/PENDIENTES.md | head -30}
      \item Más detalle: \texttt{cat /root/PENDIENTES.json}
    \end{itemize}
  \item \textbf{Quiero que Claude recuerde algo}
    \begin{itemize}
      \item ``recuerda que [X]'' en el chat → guardado en \texttt{\textasciitilde/.claude/projects/.../memory/}
    \end{itemize}
  \item \textbf{Error misterioso}
    \begin{itemize}
      \item \texttt{make logs} → \texttt{docker compose ps} → pegar error en chat
    \end{itemize}
  \item \textbf{Cuánto estoy gastando}
    \begin{itemize}
      \item ``cuánto costó esta sesión'' en el chat
    \end{itemize}
  \item \textbf{Terminar sesión bien}
    \begin{itemize}
      \item ``haz clock-out'' en el chat
    \end{itemize}
  \item \textbf{Harness desactualizado}
    \begin{itemize}
      \item \texttt{/evolve} (mejor al final de sesión larga)
    \end{itemize}
\end{itemize}

\section{🧠 Tips ADHD+AACC}

\subsection*{1. Voz → Acción Sin Fricción}

En el chat, usa el prefijo \texttt{!} para ejecutar comandos directamente:

\begin{lstlisting}
! make doctor
! cat /root/PENDIENTES.md
! docker compose ps
\end{lstlisting}

Dicta el comando y el output aparece en la conversación.

\subsection*{2. Una Tarea Activa — El Sistema Te Protege}

Cuando surge idea secundaria:
\begin{enumerate}
  \item Di: ``anota esto para después: [idea]''
  \item Claude la registra en PENDIENTES.md
  \item Tú continúas con lo que estabas haciendo
\end{enumerate}

PENDIENTES.md es tu buffer externo de memoria.

\subsection*{3. Recuperar Contexto Entre Sesiones en 30 Segundos}

\begin{lstlisting}
cat /root/SESSION_HANDOFF.md
\end{lstlisting}

El archivo contiene literalmente: ``el siguiente paso es: [instrucción concreta]''. No necesitas reconstruir nada.

\subsection*{4. Señales de que Claude Entiende vs. Adivina}

\begin{tabular}{|p{3.5cm}|p{3.5cm}|}
\hline
\textbf{Claude entiende} & \textbf{Claude adivina} \\
\hline
Cita archivos específicos con rutas & Habla en genérico (``deberías...'') \\
Ejecuta comandos y muestra output real & Dice ``probablemente funciona'' \\
Propone verificación concreta & No propone forma de verificar \\
Pide confirmación antes de cambios & Actúa sin preguntar en cambios riesgosos \\
\hline
\end{tabular}

Si ves señales de ``adivina'' → pide: ``¿leíste el archivo? muéstrame qué encontraste exactamente''.

\subsection*{5. Opus vs. Sonnet — Cuándo Importa}

\begin{tabular}{|p{3cm}|p{2cm}|p{3.5cm}|}
\hline
\textbf{Situación} & \textbf{Modelo} & \textbf{Por qué} \\
\hline
Trabajo diario, cambios de código & Sonnet (default) & Suficiente, mucho más barato \\
Orquestación multi-agente compleja & Opus 4.7 & Mejor descomposición de tareas grandes \\
Documentos importantes, arquitectura & Opus 4.7 & Razonamiento notablemente superior \\
\hline
\end{tabular}

Para cambiar: escribe \texttt{/model} en el chat.

\subsection*{6. Cuando Claude Se Atasca}

Si después de 2 intentos el resultado no está bien:

\begin{enumerate}
  \item Escribe: ``Para. Dime exactamente qué estás intentando hacer y qué encontraste.''
  \item Claude reporta su estado → identificas si es comprensión o ejecución
  \item Si sigue atascado (3+ intentos) → escala a escritura directa o sub-agente
\end{enumerate}

\subsection*{7. El Chat Es Tu Terminal Principal}

No abras terminales separadas. En el chat puedes:
\begin{itemize}
  \item Pedir que ejecute comandos
  \item Ver outputs directamente
  \item Pedir que interprete resultados
  \item Pedir que tome acción basada en lo que vio
\end{itemize}

\section{🌐 Servicios y Puertos — Referencia Rápida}

\begin{center}
\small
\begin{tabular}{|l|l|r|l|}
\hline
\textbf{Servicio} & \textbf{URL Pública} & \textbf{Puerto} & \textbf{Estado Esperado} \\
\hline
Hermes Agent & hermes.el80.space & 8080 & \texttt{Up} + \texttt{/health} 200 \\
LiteLLM Router & litellm.el80.space & 4000 & \texttt{Up} + \texttt{/health} 200 \\
Website & docs.el80.space & 3001 & \texttt{Up} \\
Grafana & grafana.el80.space & 3000 & \texttt{Up} \\
FileBrowser & files.el80.space & 8095 & \texttt{Up} \\
Whisper STT & — (interno) & 9000 & \texttt{Up} \\
Prometheus & — (interno) & 9090 & \texttt{Up} \\
\hline
\end{tabular}
\end{center}

\vspace{0.3cm}

Verificar todo:
\begin{lstlisting}
make health-check   # ver estado rápido
make status         # docker compose ps
\end{lstlisting}

\section{📋 Comandos Make — Cheat Sheet}

\begin{lstlisting}
make doctor        # Diagnóstico completo (6 pasos)
make check         # Verificación rápida: build + health + lint
make health-check  # Solo health
make build-check   # Solo build Next.js
make lint-check    # Solo lint Python
make status        # docker compose ps
make logs          # Últimas 40 líneas de cada servicio
make clean-tmp     # Limpiar /tmp/claude_progress y /tmp/*.md
\end{lstlisting}

\section{📚 Glosario Mínimo}

\begin{itemize}
  \item \textbf{harness:} Conjunto de archivos que hacen que Claude trabaje bien: CLAUDE.md, AGENTS.md, PENDIENTES.json, Makefile, hooks

  \item \textbf{skill:} Capacidad especializada que se activa con frases clave (``genera tests''). No necesitas saber el nombre técnico

  \item \textbf{swarm:} Varios agentes trabajando en paralelo en subtareas. Claude lo orquesta automáticamente

  \item \textbf{sub-agente:} Agente que Claude lanza internamente. Tú no los controlas directamente

  \item \textbf{PENDIENTES.json:} Backlog estructurado con id, prioridad, estado, verificación, evidencia. Claude lo actualiza

  \item \textbf{SESSION\_HANDOFF.md:} Nota de Claude a Claude: ``lo que quedó pendiente y el siguiente paso concreto''. Recuperación en 30 segundos

  \item \textbf{Definition of Done:} 4 criterios para considerar una tarea terminada (implementación + verificación + evidencia + stack Up)

  \item \textbf{clock-in / clock-out:} Checklist de inicio y cierre de sesión. Garantiza estado conocido

  \item \textbf{MCP:} Protocolo que permite a Claude usar herramientas externas (Google Drive, Gmail, browser, DB). Ya configurado
\end{itemize}

\vspace{1cm}

\noindent\hrulefill

\noindent\small\textit{Última actualización: 25 de mayo de 2026 · Hermes Stack · erikjosehrnndz-crypto}

\end{document}
